\documentclass{article}
\usepackage[utf8]{inputenc}
\usepackage{amsmath,amssymb}
\usepackage[most]{tcolorbox}
\usepackage{geometry}
\geometry{margin=2.5cm}

\newcommand{\step}[1]{\par\noindent\hangindent=3.5em\hangafter=1%
  \makebox[3.5em][l]{\textbf{#1.}}}
\newcommand{\steptag}[1]{\hfill{\small\textsf{[#1]}}}

\newtcolorbox{proofbox}[1]{
  colback=gray!5, colframe=gray!60,
  fonttitle=\bfseries, title={#1},
  breakable, enhanced,
  left=4pt, right=4pt, top=4pt, bottom=4pt,
  before skip=10pt, after skip=10pt
}

\begin{document}

\begin{proofbox}{The projected product on A=Im P satisfies the approximate...}
\step{1} The projected product on A=Im P satisfies the approximate Jordan identity ||((a•a)•b)•a - (a•a)•(b•a)|| <= C sqrt(eta) ||a||\textasciicircum{}3 ||b||, by exact-ambient-Jordan-identity cancellation of the leading term with O(sqrt(eta)) one-hole error terms. \steptag{validated} \\[6pt]
\step{1.1} LEFT REDUCTION. Fix a,b in A=Im P (def near-fixed-algebra). Set p:=P(a\textasciicircum{}2)=a\textasciicircum{}2-h\_\{a,a\} in A, where h\_\{a,a\}=a\textasciicircum{}2-P(a\textasciicircum{}2) in Ker P (def square-hole; P(h\_\{a,a\})=0 by P\textasciicircum{}2=P, GT-Pprops). Crude bound ||p||=||P(a\textasciicircum{}2)||<=||P|| ||a\textasciicircum{}2||<=(1+C eta)||a||\textasciicircum{}2<=C||a||\textasciicircum{}2 (GT-Pprops, GT-jb-submult ||a\textasciicircum{}2||<=||a||\textasciicircum{}2). CLAIM: ((a•a)•b)•a = P((a\textasciicircum{}2∘b)∘a) + E\_L with ||E\_L|| <= C sqrt(eta)||a||\textasciicircum{}3||b||. Hence ((a•a)•b)•a = P((a\textasciicircum{}2∘b)∘a) + O(sqrt(eta)||a||\textasciicircum{}3||b||). PROOF: a•a=P(a∘a)=P(a\textasciicircum{}2)=p, so (a•a)•b=P(p∘b). Then ((a•a)•b)•a=P(P(p∘b)∘a). Write l:=h\_\{p,b\}=p∘b-P(p∘b) in Ker P (def square-hole). By P real-linear (GT-Pprops), P(P(p∘b)∘a)=P((p∘b)∘a)-P(l∘a). Using p=a\textasciicircum{}2-h\_\{a,a\} and bilinearity of ∘ (def jordan-product), (p∘b)∘a=(a\textasciicircum{}2∘b)∘a-(h\_\{a,a\}∘b)∘a, so P((p∘b)∘a)=P((a\textasciicircum{}2∘b)∘a)-P((h\_\{a,a\}∘b)∘a). Thus E\_L=-P(l∘a)-P((h\_\{a,a\}∘b)∘a). Node 1.1.1 bounds ||P(l∘a)||<=C sqrt(eta)||a||\textasciicircum{}3||b||; node 1.1.2 bounds ||P((h\_\{a,a\}∘b)∘a)||<=C sqrt(eta)||a||\textasciicircum{}3||b||. Sum: ||E\_L||<=C sqrt(eta)||a||\textasciicircum{}3||b||. Uses node 1.1.1, node 1.1.2, GT-Pprops, GT-jb-submult, def near-fixed-algebra, def square-hole, def jordan-product. \steptag{validated} \\[6pt]
\step{1.1.1} E1 (left, l-term). l:=h\_\{p,b\}=p∘b-P(p∘b) where p=P(a\textasciicircum{}2) in A=Im P and b in A (def square-hole, def near-fixed-algebra). (i) l in Ker P: P(l)=P(p∘b)-P\textasciicircum{}2(p∘b)=P(p∘b)-P(p∘b)=0 since P\textasciicircum{}2=P (GT-Pprops). (ii) Norm: ||l||<=||p∘b||+||P(p∘b)||<=||p∘b||+(1+C eta)||p∘b||<=C||p∘b||<=C||p||||b||<=C||a||\textasciicircum{}2||b|| (triangle inequality; GT-Pprops ||P||<=1+C eta; GT-jb-submult ||p∘b||<=||p||||b||; node-1.1 bound ||p||<=C||a||\textasciicircum{}2). (iii) Apply GT-FI (first insertion) with n=l in Ker P and c=a in A: ||P(l∘a)||<=C sqrt(eta)||l||||a||<=C sqrt(eta)·C||a||\textasciicircum{}2||b||·||a||=C sqrt(eta)||a||\textasciicircum{}3||b||. CONCLUSION: ||P(l∘a)||<=C sqrt(eta)||a||\textasciicircum{}3||b||. Uses GT-FI, GT-Pprops, GT-jb-submult, def square-hole, def near-fixed-algebra. \steptag{validated} \\[4pt]
\step{1.1.2} E2 (left, h\_\{a,a\}-term). The product hole h\_\{a,a\}=a∘a-P(a∘a)=a\textasciicircum{}2-P(a\textasciicircum{}2) in Ker P (def square-hole). Apply GT-onehole-51 (eq 5.1: ||P((h\_\{r,s\}∘t)∘u)||<=C sqrt(eta)||r||||s||||t||||u||) with r=s=a, t=b, u=a, so h\_\{r,s\}=h\_\{a,a\}: ||P((h\_\{a,a\}∘b)∘a)||<=C sqrt(eta)||a||·||a||·||b||·||a||=C sqrt(eta)||a||\textasciicircum{}3||b||. CONCLUSION: ||P((h\_\{a,a\}∘b)∘a)||<=C sqrt(eta)||a||\textasciicircum{}3||b||. Uses GT-onehole-51, def square-hole. \steptag{validated} \\[4pt]
\step{1.2} RIGHT REDUCTION. With a,b,p as in node 1.1, set q:=P(b∘a)=b∘a-h\_\{b,a\} in A, h\_\{b,a\}=b∘a-P(b∘a) in Ker P (def square-hole). Crude bounds ||q||=||P(b∘a)||<=C||a||||b|| and ||h\_\{b,a\}||=||b∘a-P(b∘a)||<=||b∘a||+||P(b∘a)||<=C||a||||b|| (GT-Pprops, GT-jb-submult ||b∘a||<=||a||||b||). CLAIM: (a•a)•(b•a) = P(a\textasciicircum{}2∘(b∘a)) + E\_R with ||E\_R|| <= C sqrt(eta)||a||\textasciicircum{}3||b||. PROOF: a•a=p, b•a=P(b∘a)=q, so (a•a)•(b•a)=P(p∘q). Substitute p=a\textasciicircum{}2-h\_\{a,a\}, q=b∘a-h\_\{b,a\} and expand by bilinearity (def jordan-product): p∘q=a\textasciicircum{}2∘(b∘a)-h\_\{a,a\}∘(b∘a)-a\textasciicircum{}2∘h\_\{b,a\}+h\_\{a,a\}∘h\_\{b,a\}. Apply P (real-linear, GT-Pprops): P(p∘q)=P(a\textasciicircum{}2∘(b∘a))-P(h\_\{a,a\}∘(b∘a))-P(a\textasciicircum{}2∘h\_\{b,a\})+P(h\_\{a,a\}∘h\_\{b,a\}). So E\_R=-P(h\_\{a,a\}∘(b∘a))-P(a\textasciicircum{}2∘h\_\{b,a\})+P(h\_\{a,a\}∘h\_\{b,a\}). Node 1.2.1 bounds ||P(h\_\{a,a\}∘(b∘a))||<=C sqrt(eta)||a||\textasciicircum{}3||b||. For P(a\textasciicircum{}2∘h\_\{b,a\}) decompose a\textasciicircum{}2=p+h\_\{a,a\}: P(a\textasciicircum{}2∘h\_\{b,a\})=P(p∘h\_\{b,a\})+P(h\_\{a,a\}∘h\_\{b,a\}); node 1.2.2 bounds ||P(p∘h\_\{b,a\})||<=C sqrt(eta)||a||\textasciicircum{}3||b||. Node 1.2.3 bounds ||P(h\_\{a,a\}∘h\_\{b,a\})||<=C eta||a||\textasciicircum{}3||b||<=C sqrt(eta)||a||\textasciicircum{}3||b|| (eta<=1), and this term appears twice. Summing: ||E\_R||<=C sqrt(eta)||a||\textasciicircum{}3||b||. Uses node 1.2.1, node 1.2.2, node 1.2.3, GT-Pprops, GT-jb-submult, def near-fixed-algebra, def square-hole, def jordan-product. \steptag{validated} \\[6pt]
\step{1.2.1} E3 (right, h\_\{a,a\}∘(b∘a) term). h\_\{a,a\}=a\textasciicircum{}2-P(a\textasciicircum{}2) in Ker P (def square-hole). Apply GT-onehole-52 (eq 5.2: ||P(h\_\{r,s\}∘(t∘u))||<=C sqrt(eta)||r||||s||||t||||u||) with r=s=a, t=b, u=a, so h\_\{r,s\}=h\_\{a,a\}: ||P(h\_\{a,a\}∘(b∘a))||<=C sqrt(eta)||a||·||a||·||b||·||a||=C sqrt(eta)||a||\textasciicircum{}3||b||. CONCLUSION: ||P(h\_\{a,a\}∘(b∘a))||<=C sqrt(eta)||a||\textasciicircum{}3||b||. Uses GT-onehole-52, def square-hole. \steptag{validated} \\[4pt]
\step{1.2.2} E4 (right, p∘h\_\{b,a\} term). p=P(a\textasciicircum{}2) in A=Im P with ||p||<=C||a||\textasciicircum{}2 (node 1.1); h\_\{b,a\}=b∘a-P(b∘a) in Ker P with ||h\_\{b,a\}||<=C||a||||b|| (node 1.2, def square-hole). Since ∘ is commutative (def jordan-product), p∘h\_\{b,a\}=h\_\{b,a\}∘p. Apply GT-FI (first insertion) with n=h\_\{b,a\} in Ker P and c=p in A: ||P(h\_\{b,a\}∘p)||<=C sqrt(eta)||h\_\{b,a\}||||p||<=C sqrt(eta)·C||a||||b||·C||a||\textasciicircum{}2=C sqrt(eta)||a||\textasciicircum{}3||b||. Hence ||P(p∘h\_\{b,a\})||<=C sqrt(eta)||a||\textasciicircum{}3||b||. Uses GT-FI, def square-hole, def near-fixed-algebra, def jordan-product. \steptag{validated} \\[4pt]
\step{1.2.3} E5 (right, two-hole term h\_\{a,a\}∘h\_\{b,a\}). h\_\{a,a\}=a\textasciicircum{}2-P(a\textasciicircum{}2)=-q\_\{a,a\} and h\_\{b,a\}=b∘a-P(b∘a)=-q\_\{b,a\}, where q\_\{r,s\}=P(r∘s)-r∘s (def square-hole). So h\_\{a,a\}∘h\_\{b,a\}=q\_\{a,a\}∘q\_\{b,a\}. Apply GT-HH (two-hole) with (r,s)=(a,a) and (u,v)=(b,a): ||P(q\_\{a,a\}∘q\_\{b,a\})||<=C eta||a||·||a||·||b||·||a||=C eta||a||\textasciicircum{}3||b||. Since eta<=eta\_0<=1 we have eta<=sqrt(eta), so ||P(h\_\{a,a\}∘h\_\{b,a\})||<=C eta||a||\textasciicircum{}3||b||<=C sqrt(eta)||a||\textasciicircum{}3||b||. CONCLUSION: ||P(h\_\{a,a\}∘h\_\{b,a\})||<=C eta||a||\textasciicircum{}3||b||<=C sqrt(eta)||a||\textasciicircum{}3||b||. Uses GT-HH, def square-hole. \steptag{validated} \\[4pt]
\step{1.3} CANCELLATION + ASSEMBLY. CLAIM: P((a\textasciicircum{}2∘b)∘a)=P(a\textasciicircum{}2∘(b∘a)) EXACTLY, hence ||((a•a)•b)•a-(a•a)•(b•a)|| <= C sqrt(eta)||a||\textasciicircum{}3||b||. PROOF of cancellation: node 1.3.1 proves the symmetric ambient Jordan identity (a\textasciicircum{}2∘b)∘a=a\textasciicircum{}2∘(b∘a) as an exact identity of elements of B(H)\_sa (derived from GT-jordan-identity + commutativity GT-jordan-product). Applying the well-defined real-linear map P to both sides of this exact equality of elements of B(H)\_sa gives P((a\textasciicircum{}2∘b)∘a)=P(a\textasciicircum{}2∘(b∘a)). ASSEMBLY: by node 1.1, ((a•a)•b)•a=P((a\textasciicircum{}2∘b)∘a)+E\_L; by node 1.2, (a•a)•(b•a)=P(a\textasciicircum{}2∘(b∘a))+E\_R; with ||E\_L||,||E\_R||<=C sqrt(eta)||a||\textasciicircum{}3||b||. Subtracting and using the cancellation P((a\textasciicircum{}2∘b)∘a)=P(a\textasciicircum{}2∘(b∘a)): ((a•a)•b)•a-(a•a)•(b•a)=E\_L-E\_R, so ||((a•a)•b)•a-(a•a)•(b•a)||<=||E\_L||+||E\_R||<=C sqrt(eta)||a||\textasciicircum{}3||b||. This is exactly the contract (JB4, def eps-jb-algebra). Uses node 1.1, node 1.2, node 1.3.1, GT-jordan-identity, GT-jordan-product, def near-fixed-algebra, def jordan-product. \steptag{validated} \\[6pt]
\step{1.3.1} SYMMETRIC AMBIENT JORDAN IDENTITY (exact). CLAIM: in B(H)\_sa with the ambient Jordan product ∘ (def jordan-product), (a\textasciicircum{}2∘b)∘a=a\textasciicircum{}2∘(b∘a) for all a,b, where a\textasciicircum{}2=a∘a. PROOF: GT-jordan-identity is HOS eq (2.18): a∘(b∘a\textasciicircum{}2)=(a∘b)∘a\textasciicircum{}2 (an exact identity valid for the special Jordan product on any associative algebra, in particular B(H)). The product ∘ is commutative: a∘b=b∘a (GT-jordan-product, since (1/2)(ab+ba)=(1/2)(ba+ab)). Apply commutativity to BOTH sides of (2.18): LHS a∘(b∘a\textasciicircum{}2)=(b∘a\textasciicircum{}2)∘a=(a\textasciicircum{}2∘b)∘a (inner b∘a\textasciicircum{}2=a\textasciicircum{}2∘b by commutativity, then outer a∘X=X∘a); RHS (a∘b)∘a\textasciicircum{}2=a\textasciicircum{}2∘(a∘b)=a\textasciicircum{}2∘(b∘a) (outer commutativity, then inner a∘b=b∘a). Therefore (a\textasciicircum{}2∘b)∘a=a\textasciicircum{}2∘(b∘a), an EXACT equality of elements of B(H)\_sa. (This is one commutativity-relabeling step from (2.18), no error term.) Uses GT-jordan-identity, GT-jordan-product, def jordan-product. \steptag{validated} \\[4pt]
\end{proofbox}

\end{document}
